% GENERATED FILE. Edit canonical lessons, then run npm run generate:books.
% canonical-source-hash: fnv1a64:7f179ed0664d688f
% canonical-lessons: GU-C200-tapali, GU-C200-parbidiyu, GU-C200-parsal, GU-C200-nakhvu, GU-C200-nondhavvu, GU-R200-review-195-200

\chapter{At the post office}
\label{ch:gu-at-the-post-office}
\hlchaptermodality{200}

\begin{chapteropening}
\textbf{By the end of this chapter:} \emph{I can post a parcel and register a letter.}

Complete the last lesson of chapter 200: I can post a parcel and register a letter.
\end{chapteropening}

\section[\texorpdfstring{\gu{ટપાલી} (ṭapālī)}{ટપાલી (ṭapālī)}]{\gu{ટપાલી} (ṭapālī) — postman}

\label{lesson:GU-C200-tapali}



\begin{quote}
Before the new one: say the Gujarati for to sneak in, then the Gujarati for to move aside.
\end{quote}

\subsection*{You'll want to know: \gu{ટપાલી}}
\textbf{\gu{ટપાલી}} — \emph{ṭapālī} — \textquotedblleft{}postman\textquotedblright{}.

\begin{quote}
\textbf{\gu{ટપાલી} \gu{સવારે} \gu{ટપાલ} \gu{લાવે} \gu{છે}.} — \emph{ṭapālī savāre ṭapāl lāve che} — \textquotedblleft{}The postman brings the post in the morning.\textquotedblright{}
\end{quote}

Built on \textbf{\gu{ટપાલ}}, \emph{ṭapāl}, 'post'.

\begin{cousinweb}
\textbf{\gu{ટપાલી}} (\emph{ṭapālī}) is built on \emph{ṭapāl}, \textquotedblleft{}post, mail\textquotedblright{}, with the agent ending \emph{-ī}. The origin of \emph{ṭapāl} is not settled: neither CDIAL nor DEDR has an entry for it.
\end{cousinweb}

\begin{grammarlens}[title={What you've built}]
One more word for everyday life at A2.
\end{grammarlens}

\subsection*{Your turn}
\begin{itemize}
\raggedright
  \item \emph{Say it:} \emph{ṭapālī}
  \item \emph{Say it:} \emph{ṭapālī}, once more
  \item \emph{Say it:} \emph{ṭapālī savāre ṭapāl lāve che} — The postman brings the post in the morning.
  \item \emph{Recall:} say the Gujarati for to sneak in, then the Gujarati for to move aside, then say \emph{ṭapālī} again
\end{itemize}

\begin{tcolorbox}[breakable,colback=teal!4,colframe=teal!35!black,title={Before you move on}]
What does \gu{ટપાલી} mean? (\textquotedblleft{}Postman\textquotedblright{}.) Say it once more.
\end{tcolorbox}
\section[\texorpdfstring{\gu{પરબીડિયું} (parbīḍiyũ)}{પરબીડિયું (parbīḍiyũ)}]{\gu{પરબીડિયું} (parbīḍiyũ) — envelope}

\label{lesson:GU-C200-parbidiyu}



\begin{quote}
Before the new one: say the Gujarati for to move aside, then the Gujarati for postman.
\end{quote}

\subsection*{You'll want to know: \gu{પરબીડિયું}}
\textbf{\gu{પરબીડિયું}} — \emph{parbīḍiyũ} — \textquotedblleft{}envelope\textquotedblright{}.

\begin{quote}
\textbf{\gu{પત્ર} \gu{પરબીડિયામાં} \gu{છે}.} — \emph{patra parbīḍiyāmā̃ che} — \textquotedblleft{}The letter is in the envelope.\textquotedblright{}
\end{quote}

A neuter noun.

\begin{cousinweb}
The origin of \textbf{\gu{પરબીડિયું}} (\emph{parbīḍiyũ}) is not settled. CDIAL has no entry for it, and EtymDB lists no source.
\end{cousinweb}

\begin{grammarlens}[title={What you've built}]
One more word for everyday life at A2.
\end{grammarlens}

\subsection*{Your turn}
\begin{itemize}
\raggedright
  \item \emph{Say it:} \emph{parbīḍiyũ}
  \item \emph{Say it:} \emph{parbīḍiyũ}, once more
  \item \emph{Say it:} \emph{patra parbīḍiyāmā̃ che} — The letter is in the envelope.
  \item \emph{Recall:} say the Gujarati for to move aside, then the Gujarati for postman, then say \emph{parbīḍiyũ} again
\end{itemize}

\begin{tcolorbox}[breakable,colback=teal!4,colframe=teal!35!black,title={Before you move on}]
What does \gu{પરબીડિયું} mean? (\textquotedblleft{}Envelope\textquotedblright{}.) Say it once more.
\end{tcolorbox}
\section[\texorpdfstring{\gu{પાર્સલ} (pārsal)}{પાર્સલ (pārsal)}]{\gu{પાર્સલ} (pārsal) — parcel}

\label{lesson:GU-C200-parsal}



\begin{quote}
Before the new one: say the Gujarati for postman, then the Gujarati for envelope.
\end{quote}

\subsection*{You'll want to know: \gu{પાર્સલ}}
\textbf{\gu{પાર્સલ}} — \emph{pārsal} — \textquotedblleft{}parcel\textquotedblright{}.

\begin{quote}
\textbf{\gu{મેં} \gu{ગઈકાલે} \gu{પાર્સલ} \gu{મોકલ્યું}.} — \emph{mẽ gaīkāle pārsal mokalyũ} — \textquotedblleft{}I sent a parcel yesterday.\textquotedblright{}
\end{quote}

\begin{cousinweb}
\textbf{\gu{પાર્સલ}} (\emph{pārsal}) is a loan from English \emph{parcel}.
\end{cousinweb}

\begin{grammarlens}[title={What you've built}]
One more word for everyday life at A2.
\end{grammarlens}

\subsection*{Your turn}
\begin{itemize}
\raggedright
  \item \emph{Say it:} \emph{pārsal}
  \item \emph{Say it:} \emph{pārsal}, once more
  \item \emph{Say it:} \emph{mẽ gaīkāle pārsal mokalyũ} — I sent a parcel yesterday.
  \item \emph{Recall:} say the Gujarati for postman, then the Gujarati for envelope, then say \emph{pārsal} again
\end{itemize}

\begin{tcolorbox}[breakable,colback=teal!4,colframe=teal!35!black,title={Before you move on}]
What does \gu{પાર્સલ} mean? (\textquotedblleft{}Parcel\textquotedblright{}.) Say it once more.
\end{tcolorbox}
\section[\texorpdfstring{\gu{નાખવું} (nākhvũ)}{નાખવું (nākhvũ)}]{\gu{નાખવું} (nākhvũ) — to put in, to drop in}

\label{lesson:GU-C200-nakhvu}



\begin{quote}
Before the new one: say the Gujarati for envelope, then the Gujarati for parcel.
\end{quote}

\subsection*{You'll want to know: \gu{નાખવું}}
\textbf{\gu{નાખવું}} — \emph{nākhvũ} — \textquotedblleft{}to put in, to drop in\textquotedblright{}.

\begin{quote}
\textbf{\gu{તે} \gu{પત્ર} \gu{ટપાલપેટીમાં} \gu{નાખે} \gu{છે}.} — \emph{te patra ṭapālpeṭīmā̃ nākhe che} — \textquotedblleft{}She drops the letter into the postbox.\textquotedblright{}
\end{quote}

\begin{cousinweb}
\textbf{\gu{નાખવું}} (\emph{nākhvũ}) is a tadbhava. CDIAL lists it, glossed \textquotedblleft{}to throw, put on, leave\textquotedblright{}, under Sanskrit \emph{naṅkṣyati}, \textquotedblleft{}will perish\textquotedblright{} (CDIAL 6930).
\end{cousinweb}

\begin{grammarlens}[title={What you've built}]
One more word for everyday life at A2.
\end{grammarlens}

\subsection*{Your turn}
\begin{itemize}
\raggedright
  \item \emph{Say it:} \emph{nākhvũ}
  \item \emph{Say it:} \emph{nākhvũ}, once more
  \item \emph{Say it:} \emph{te patra ṭapālpeṭīmā̃ nākhe che} — She drops the letter into the postbox.
  \item \emph{Recall:} say the Gujarati for envelope, then the Gujarati for parcel, then say \emph{nākhvũ} again
\end{itemize}

\begin{tcolorbox}[breakable,colback=teal!4,colframe=teal!35!black,title={Before you move on}]
What does \gu{નાખવું} mean? (\textquotedblleft{}To put in, to drop in\textquotedblright{}.) Say it once more.
\end{tcolorbox}
\section[\texorpdfstring{\gu{નોંધાવવું} (nõdhāvvũ)}{નોંધાવવું (nõdhāvvũ)}]{\gu{નોંધાવવું} (nõdhāvvũ) — to register, to book}

\label{lesson:GU-C200-nondhavvu}



\begin{quote}
Before the new one: say the Gujarati for parcel, then the Gujarati for to put in.
\end{quote}

\subsection*{You'll want to know: \gu{નોંધાવવું}}
\textbf{\gu{નોંધાવવું}} — \emph{nõdhāvvũ} — \textquotedblleft{}to register, to book\textquotedblright{}.

\begin{quote}
\textbf{\gu{મેં} \gu{પત્ર} \gu{નોંધાવ્યો}.} — \emph{mẽ patra nõdhāvyo} — \textquotedblleft{}I had the letter registered.\textquotedblright{}
\end{quote}

\begin{cousinweb}
\textbf{\gu{નોંધાવવું}} (\emph{nõdhāvvũ}) is the causative of \emph{nõdhvũ}, \textquotedblleft{}to note down, enter in a book\textquotedblright{}, which CDIAL lists under \emph{*nibandhati}, beside Sanskrit \emph{níbadhnāti}, \textquotedblleft{}binds on; writes down\textquotedblright{} (CDIAL 7223).
\end{cousinweb}

\begin{grammarlens}[title={What you've built}]
One more word for everyday life at A2.
\end{grammarlens}

\subsection*{Your turn}
\begin{itemize}
\raggedright
  \item \emph{Say it:} \emph{nõdhāvvũ}
  \item \emph{Say it:} \emph{nõdhāvvũ}, once more
  \item \emph{Say it:} \emph{mẽ patra nõdhāvyo} — I had the letter registered.
  \item \emph{Recall:} say the Gujarati for parcel, then the Gujarati for to put in, then say \emph{nõdhāvvũ} again
\end{itemize}

\begin{tcolorbox}[breakable,colback=teal!4,colframe=teal!35!black,title={Before you move on}]
What does \gu{નોંધાવવું} mean? (\textquotedblleft{}To register, to book\textquotedblright{}.) Say it once more.
\end{tcolorbox}
\section[(dialogue)]{Words from chapters 195–200 — a review}

\label{lesson:GU-R200-review-195-200}



\begin{quote}
Say the Gujarati for to get lost and to roam.
\end{quote}

\subsection*{Your turn}
Say these aloud:

\begin{itemize}
\raggedright
  \item to get lost, to roam, to show, whereabouts, distance
  \item building, storey, elevator, to go in, entry
  \item watchman, to check, identity card, closed, to open up
  \item signboard, announcement, advertisement, ban, danger
  \item crowd, queue, rule, to sneak in, to move aside
  \item postman, envelope, parcel, to put in, to register
\end{itemize}

\begin{tcolorbox}[breakable,colback=teal!4,colframe=teal!35!black,title={Before you move on}]
To get lost? (\textbf{\gu{ખોવાવું}}.) To roam? (\textbf{\gu{રખડવું}}.) To show? (\textbf{\gu{બતાવવું}}.) Whereabouts? (\textbf{\gu{ઠેકાણું}}.) Distance? (\textbf{\gu{અંતર}}.) Building? (\textbf{\gu{મકાન}}.) Storey? (\textbf{\gu{માળ}}.) Elevator? (\textbf{\gu{લિફ્ટ}}.) To go in? (\textbf{\gu{પેસવું}}.) Entry? (\textbf{\gu{પ્રવેશ}}.) Watchman? (\textbf{\gu{ચોકીદાર}}.) To check? (\textbf{\gu{તપાસવું}}.) Identity card? (\textbf{\gu{ઓળખપત્ર}}.) Closed? (\textbf{\gu{બંધ}}.) To open up? (\textbf{\gu{ખૂલવું}}.) Signboard? (\textbf{\gu{પાટિયું}}.) Announcement? (\textbf{\gu{સૂચના}}.) Advertisement? (\textbf{\gu{જાહેરાત}}.) Ban? (\textbf{\gu{મનાઈ}}.) Danger? (\textbf{\gu{જોખમ}}.) Crowd? (\textbf{\gu{ભીડ}}.) Queue? (\textbf{\gu{કતાર}}.) Rule? (\textbf{\gu{નિયમ}}.) To sneak in? (\textbf{\gu{ઘૂસવું}}.) To move aside? (\textbf{\gu{ખસવું}}.) Postman? (\textbf{\gu{ટપાલી}}.) Envelope? (\textbf{\gu{પરબીડિયું}}.) Parcel? (\textbf{\gu{પાર્સલ}}.) To put in? (\textbf{\gu{નાખવું}}.) To register? (\textbf{\gu{નોંધાવવું}}.)
\end{tcolorbox}
